\documentclass[aspectratio=169,10pt]{beamer}
\usetheme[numbering=fraction,progressbar=frametitle]{metropolis}
\usepackage{fontspec}
\setsansfont{Latin Modern Sans}
\usepackage{amsmath,amssymb,mathtools,graphicx,booktabs}
\usepackage{polyglossia}
\setdefaultlanguage{french}
\setbeamertemplate{navigation symbols}{}
\setbeamerfont{frametitle}{size=\large}
\setbeamerfont{normal text}{size=\normalsize}
\setlength{\parskip}{.45em}
\hypersetup{pdfauthor={},pdfsubject={Lecture slides},pdfkeywords={},pdftitle={Représentations paramétriques et équations cartésiennes}}
\providecommand{\N}{\mathbb{N}}
\providecommand{\Z}{\mathbb{Z}}
\providecommand{\D}{\mathbb{D}}
\providecommand{\Q}{\mathbb{Q}}
\providecommand{\R}{\mathbb{R}}
\providecommand{\C}{\mathbb{C}}
\providecommand{\vect}[1]{\overrightarrow{#1}}
\providecommand{\vi}{\vec{\imath}}
\providecommand{\vj}{\vec{\jmath}}
\providecommand{\vk}{\vec{k}}
\providecommand{\norm}[1]{\left\lVert #1\right\rVert}
\providecommand{\vc}[2]{\begin{pmatrix}#1\\#2\end{pmatrix}}
\providecommand{\deter}[2]{\det\left(#1,#2\right)}
\providecommand{\abs}[1]{\left\lvert #1\right\rvert}
\providecommand{\dd}{\,\mathrm{d}}
\providecommand{\e}{\mathrm{e}}
\providecommand{\Card}{\operatorname{Card}}
\providecommand{\E}{\operatorname{E}}
\providecommand{\V}{\operatorname{V}}
\providecommand{\degre}{^{\circ}}
\providecommand{\un}[1]{\,\mathrm{#1}}
\providecommand{\sci}[2]{#1\times10^{#2}}
\providecommand{\pgcd}{\operatorname{PGCD}}
\providecommand{\ppcm}{\operatorname{PPCM}}
\title{Représentations paramétriques et équations cartésiennes}
\subtitle{Cours et exemples}
\date{}
\author{}
\begin{document}
\begin{frame}[plain]
\vfill
{\usebeamerfont{title}\inserttitle\par}
\medskip
{\usebeamerfont{subtitle}\insertsubtitle\par}
\medskip
{\usebeamercolor[fg]{progress bar}\rule{\textwidth}{.5pt}}
\vfill
\end{frame}
\begin{frame}{Objectifs}
Décrire une droite et un plan ; déterminer des intersections ; choisir une méthode pour une projection.
\end{frame}
\begin{frame}{Notes de cours}
Dans tout ce chapitre, les coordonnées sont données dans un repère orthonormé de l'espace. Cette hypothèse est nécessaire pour identifier un vecteur normal à partir des coefficients d'une équation et pour calculer les distances.
\end{frame}
\begin{frame}{Droites paramétrées}
La droite passant par \(A(x_A;y_A;z_A)\) de vecteur directeur non nul \((a;b;c)\) admet la représentation
\[d:\quad\begin{cases}x=x_A+at\\y=y_A+bt\\z=z_A+ct\end{cases}\quad(t\in\R).\]
Un point appartient à la droite s'il existe \emph{un même} réel \(t\) satisfaisant les trois égalités. Une droite admet de nombreuses représentations : changement du point de départ ou multiplication du directeur par un réel non nul.
\end{frame}
\begin{frame}{Droites paramétrées}
\((x;y;z)=(1;2;0)+t(2;-1;3)\) décrit la même droite que \((3;1;3)+s(4;-2;6)\) ; le lien est \(t=1+2s\).
\end{frame}
\begin{frame}{Plans cartésiens}
Si \((a;b;c)\ne(0;0;0)\), l'ensemble \(ax+by+cz+d=0\) est un plan de vecteur normal \((a;b;c)\). Le plan de normal \(\vec n\) passant par \(A\) se construit en écrivant \(\vect{AM}\cdot\vec n=0\).
\end{frame}
\begin{frame}{Plans cartésiens}
La condition \(\vect{AM}\cdot\vec n=0\) s'écrit \(a(x-x_A)+b(y-y_A)+c(z-z_A)=0\). En développant, \(d=-ax_A-by_A-cz_A\). Réciproquement, tout point satisfaisant cette équation vérifie la condition d'orthogonalité, donc appartient au plan.
\end{frame}
\begin{frame}{Plans cartésiens}
Le plan de normal \((2;-1;3)\) passant par \(A(1;2;0)\) est \(2(x-1)-(y-2)+3z=0\), soit \(2x-y+3z=0\).
\end{frame}
\begin{frame}{Intersections}
Pour une droite et un plan, remplacer \(x,y,z\) dans l'équation du plan par leurs expressions paramétriques. Une solution en \(t\) donne une intersection ponctuelle ; une identité donne une droite incluse ; une contradiction donne une droite parallèle disjointe.
Pour deux plans, des normales non colinéaires donnent une intersection suivant une droite ; résoudre le système à deux équations en choisissant une variable comme paramètre. Si les normales sont colinéaires, les plans sont parallèles : distinguer confondus et disjoints.
Pour deux droites, utiliser des paramètres \(s,t\) indépendants et résoudre les trois équations. Sans solution et avec directeurs non colinéaires, les droites sont non coplanaires.
\end{frame}
\begin{frame}{Projection sur un plan}
Si \(\mathcal P:ax+by+cz+d=0\), chercher le projeté de \(M\) sur la droite normale passant par \(M\) : \(H=M+t(a;b;c)\). Remplacer dans l'équation du plan pour déterminer \(t\), puis calculer \(MH\).
On en déduit, sans en faire une formule à apprendre isolément,
\[d(M,\mathcal P)=\dfrac{|ax_M+by_M+cz_M+d|}{\sqrt{a^2+b^2+c^2}}.\]
\end{frame}
\begin{frame}{Projection sur un plan}
Après le calcul d'un projeté, contrôler simultanément \(H\in\mathcal P\) et la colinéarité de \(\vect{MH}\) avec le normal.
\end{frame}
\begin{frame}{Activité de recherche}
\textbf{Trois droites face à un plan - 30 min.}

Dans un repère orthonormé de l'espace, on considère le plan \(\mathcal P\) d'équation \(x+y+z=3\). Dans chaque question, le paramètre \(t\) parcourt \(\R\).
\end{frame}
\begin{frame}{Activité de recherche (2)}
\begin{enumerate}
\setcounter{enumi}{0}\item La droite \(d_1\) admet la représentation \((x;y;z)=(t;t;t)\). Déterminer \(d_1\cap\mathcal P\).

\item La droite \(d_2\) admet la représentation \((x;y;z)=(3+t;-t;0)\). Déterminer \(d_2\cap\mathcal P\).

\item La droite \(d_3\) admet la représentation \((x;y;z)=(t;-t;0)\). Déterminer \(d_3\cap\mathcal P\).

\item Pour chaque droite, préciser si l'équation obtenue après substitution dans celle du plan a une solution unique, aucune solution ou tout réel pour solution. Relier ces cas aux positions géométriques trouvées.

\end{enumerate}
\end{frame}
\end{document}
